\chapter*{\scshape Ringraziamenti}
\addcontentsline{toc}{chapter}{\scshape Ringraziamenti}

\noindent A \textbf{tutta la mia famiglia}, vicina e lontana, che mi ha sempre sostenuto e incoraggiato. A chi \textbf{ho conosciuto tardi} e mi ha fatto sentire a casa dal primo giorno.\\

\medskip
\noindent A \textbf{mio padre}, che senza saperlo ha acceso la mia passione. Grazie per non avermi mai fatto mancare sostegno, e per quel \textit{``manda il mondo capasotto``}.\\

\medskip
\noindent A \textbf{mia madre} che mi ha cresciuto e mi ha accompagnato fin dai primi passi, dai compiti di scuola a questo traguardo.  \\

\medskip
\noindent Il mio ringraziamento più grande va a \textbf{te, amore mio}. Tre anni fa c'eri al mio test d'ingresso, oggi sei qui alla mia laurea, e in mezzo c'è stato tutto: le giornate a studiare insieme, i concerti, i viaggi, i giorni belli e quelli brutti. Tu, che ti arrabbiavi ogni volta che mi distraevo e mi riportavi sullo studio quando la testa andava altrove.
Ti spaventavi davanti ai miei appunti, quaderni e lavagne pieni di formule, e io sorridevo e cercavo di spiegarti quello che stava scritto, come quando ero il tuo professore personale di matematica.\\
Eppure, nonostante tutto, \textbf{ci sei sempre per me}, a incoraggiarmi, a spronarmi, a sostenermi, a dirmi che sono bravissimo, anche quando l'ansia di non riuscirci era tanta. Sei la mia forza, la prima persona che chiamavo dopo ogni esame. \\
Grazie per avermi fatto crescere, per avermi insegnato cosa vuol dire amare e per avermi tenuto la mano anche quando non lo meritavo. \\
Sono orgoglioso di te e del coraggio con cui insegui i tuoi sogni, più di quanto riesca a scrivere qui. 
Qualunque strada sceglierai, non smettere mai di crederci, nemmeno nei giorni d'ansia: io sarò al tuo fianco, come tu sei sempre stata al mio.
Non so dove saremo tra qualche anno. So che ci saremo noi, e che questo traguardo è il primo mattone del nostro futuro.
Oggi la corona la porto io, ma non vedo l'ora di metterla a te. \textbf{Ti amo amore mio, oggi e per sempre.}\\
\medskip

\noindent A \textbf{mio fratello} coltiva la tua curiosità ogni giorno, e non perdere mai la strada. Ricorda che tutti noi ci saremo sempre.\\

\medskip
\noindent Agli \textbf{amici storici}, con cui ho passato quasi dieci anni e con cui sono cresciuto, anche se oggi ognuno sta prendendo la propria strada: le strade cambiano, le storie vissute insieme restano. \\
Ad \textbf{\textit{Andrea e a Genbino}}, compagni di studio con cui oggi tagliamo lo stesso traguardo. \\

\medskip
\noindent E \textbf{ai nuovi amici}, che in questi anni mi hanno tenuto compagnia.\\

\medskip
\noindent A \textbf{chi in questi anni mi ha dato fiducia} e occasioni per crescere, anche fuori dall'università. Grazie a tutti i \textbf{professori} che ho incontrato in questi anni, che mi hanno aiutato a crescere. \\

%\medskip
%\noindent E \textbf{a chi ha un sogno} e ha paura di non farcela: date corpo alle vostre idee, l'unico fallimento è non provarci.\\

\medskip
\noindent Questo è solo il primo passo. Ci vediamo \textbf{alla magistrale}.